% Paper 6 (GRPO-Registry) standalone introduction + related work
\section{Introduction}
\label{sec:p6-intro}
Group Relative Policy Optimization~\citep{shao2024deepseekmath, guo2025deepseekr1}
has spawned a family of variants at a pace that outstrips the community's
bookkeeping. DAPO changes the clip asymmetry, the sampling loop, the loss
aggregation, and the reward shaping in one release~\citep{yu2025dapo}; Dr.~GRPO
deletes two normalization terms~\citep{tong2025drgrpo}; GSPO moves the
importance ratio from tokens to sequences~\citep{qwen2025gspo}; and a long tail
of further variants adjusts baselines, curricula, or advantage estimators
\citep{lin2025cppo, nan2025ngrpo, zhang2025scaffgrpo, mcgrpo2025, gift2025,
liu2026gdpo, ichihara2025mogrpo, lu2025grpolead, bytedance2025vapo,
hu2026openvlthinker, le2025rlzvp}. Each delta is documented---but in prose, in
an appendix, in a YAML default, or in a code path that a framework may or may
not implement. There is no machine-readable place to ask: \emph{which changes
does this stack actually make to base GRPO, and which of them does my launcher
actually execute?}

This gap is not cosmetic. Across the \ours{} program we repeatedly measured the
same algorithm label producing different updates. Holding model, group size,
learning rate, data, and seed fixed and swapping only the backend moved final
GSM8K reward from $84.4\%$ to $5.0\%$---a $17\times$ swing under an identical
label. Running ``DAPO'' on an open Colab trainer whose dynamic sampling
demonstrably filters zero-variance groups produced mean ZVF $0.00$; running
``DAPO'' on a closed managed stack, where only an asymmetric-clip surrogate is
user-settable, produced mean ZVF $0.58$. Same label, opposite telemetry. When
labels underdetermine update rules, empirical rankings of ``GRPO vs.\ DAPO
vs.\ GSPO'' are rankings of stacks, not algorithms---a confound our companion
pillar papers (P1 scaling, P2 ZVF, P3 group size, P4 length bias) hit from four
different directions.

We respond with a \emph{resource}, not another variant:
\textsc{GRPO-Registry}, a machine-readable catalog with two record types.
A \emph{stack record} describes one deployed training stack via the seven-item
minimum reportable stack (MIN-REPORT-RL): (1) loss form, (2) reference policy
and KL handling, (3) sampler/backend/precision, (4) per-step ZVF/GU telemetry,
(5) group-size schedule, (6) held-out split disjoint from the reward
environment, and (7) decontamination plus parser-robustness probes. A
\emph{variant-delta record} pins down, component by component, what a named
variant changes relative to base GRPO, verified against the source paper; stack
records then declare per-component implementation status
(\texttt{implemented} / \texttt{surrogate} / \texttt{absent} / \texttt{unknown}),
which is exactly where label flips become visible to a script.

Our contributions are:
(i) the schema (Section~\ref{sec:p6-schema});
(ii) verified delta records for DAPO, Dr.~GRPO, and GSPO;
(iii) twelve seed entries populated from real artifacts of this benchmark---the
released per-framework config dumps, an open-trainer audit, and a
managed-runtime head-to-head (Section~\ref{sec:p6-population});
(iv) a query API, a 0--100 MIN-REPORT auditor badge, and a \texttt{stackdiff}
tool that converts a pair of entries into a flip-risk verdict on the R0--R5
ladder (Section~\ref{sec:p6-query}); and
(v) worked examples in which the registry flags, \emph{before any GPU time is
spent}, the label-flip risk that our experiments later confirmed empirically
(Section~\ref{sec:p6-worked}).

\subsection{Related Work}
\label{sec:p6-related}
Structured reporting artifacts exist for models and
datasets---model cards~\citep{mitchell2019modelcards} and
datasheets~\citep{gebru2021datasheets}---and evaluation has consolidated around
shared harnesses and living benchmarks such as HELM~\citep{liang2022helm} and
the LM Evaluation Harness~\citep{gao2021harness, biderman2024lessons}. Papers
with Code catalogued paper--code--leaderboard linkage at
scale~\citep{paperswithcode2025}. None of these describes the \emph{training
stack} of an RL fine-tuning run: what loss form, KL handling, and sampling loop
were in force. Analyses that reduce GRPO variants to a shared functional core
\citep{wu2025grpo_dpo, liu2026gdpo} are complementary: they say when two update
rules are equivalent in theory, while the registry records which rule a stack
runs in practice. Section~\ref{sec:p6-related-resources} gives a fuller
comparison.
